\chapter{Four Pressures, One Ledger}
\label{ch:stories}
% ===================================================================

Chapter~\ref{chap:pledge-ch5} of the Pledge told four stories about the origin of human intelligence, in the register they deserve on first hearing. This chapter tells them again in the register of the Turn, and it does so for a reason that Chapter~\ref{ch:trick} left on the table.

Observation~\ref{obs:soluble} said that the reference bootstrap problem has been solved at least twice in nature, and Section~\ref{sec:soluble} extracted three features common to the solutions: the solver was a population, the interface was not chosen, and there was an external criterion. It then said that what is established is solubility \emph{by populations, under selection, over long timescales, with co-evolving interfaces}, and stopped, because a list of features is not a mechanism. Chapter~\ref{ch:lie} recorded one candidate mechanism, a liar. This chapter records what the people who actually study the human case have to say about the mechanism, and reads it through the ledger.

The claim i am going to make is modest in form and, i think, consequential in content. It is that each of the three surviving stories names a \emph{pressure on the account} of a mortal scientist, that the three pressures are already present in the Turn as distinct constructions, and that the three-beat schema the journalist Max Fisher uses to tell the stories --- a pressure, a trigger, what intelligence is --- is the schema this book has been using all along under other names \cite{fisher2026}. Where that identification holds, the anthropology and the calculus are two witnesses to one shape. Where it does not hold, the mismatch is informative, and i will say where it does not.

\section{The schema is the book's}
\label{st:sec:schema}

Fisher tells each story in three beats. First an \emph{environmental pressure}: the thing the world did that made the old way of living stop paying. Then an \emph{intelligence trigger}: the thing a lineage evolved in response, which was not itself intelligence but became it. Then a conclusion about \emph{what intelligence is}: what it is optimized for, what its signature is, what it says about the animal that has it.

Read against the Turn, the beats have exact counterparts.

\begin{principle}[Pressure, trigger, class]
\label{st:prin:schema}
An environmental pressure is a constraint on the account of a mortal scientist: a change in what the funding gate admits, in what a harvest costs, or in who can debit whom. An intelligence trigger is a construction in cost-accounted rho that a population adopts under that constraint. What intelligence is, under a story, is the bisimulation class the construction induces: what an observer with the book's instruments can tell apart about a learner that has adopted it, and cannot tell apart about one that has not.
\end{principle}

\noindent
The principle is not a theorem; it is a reading rule, and the test of it is whether it makes the stories sharper rather than vaguer. The rest of the chapter is that test.

One premise comes first, because the three surviving stories share it and the schema needs it. Fisher takes from a 2025 study by Gonz\'alez-Forero and G\'omez-Robles the claim that human brain size shows no signature of direct selection: an evo-devo model that recovers the brain and body sizes of seven hominin species does so with selection acting on reproductive traits and brain size riding along as a genetically correlated byproduct, a spandrel in Gould and Lewontin's sense \cite{gonzalezforero2024,gonzalezforero2025}. The authors are careful to say that a model with direct selection on cognition might fit better and has not been compared; their earlier model attributed the expansion largely to ecological challenge with a cooperative component \cite{gonzalezforero2018}. i take the claim at the strength they state it and no stronger.

\begin{observation}[The byproduct premise is the ledger premise]
\label{st:obs:byproduct}
Nothing in this book selects for intelligence. The mortal scientist of Chapter~\ref{ch:sci} is selected for not running out; Proposition~\ref{suf:prop:derived} derives its reinforcement signal from the state of its own account and nothing else; Section~\ref{sec:kant-affordable} gathers every limit on what it can know and finds each to be a price. Whatever the book calls intelligence is a description, after the fact, of what a ledger under pressure was obliged to be able to do. ``Brain size is not an adaptation but a byproduct'' is that sentence in a paleontologist's vocabulary.
\end{observation}

\noindent
This is why the first story has to go before the others can be read. It is the one story in which intelligence is the thing selected.

\section{The monolith is the angle brackets}
\label{st:sec:monolith}

The first story, which Fisher calls the monolith after Kubrick and which Harari calls the Tree of Knowledge mutation \cite{harari2014}, has a precise academic form: a neural mutation around fifty thousand years ago that rewired the sapiens brain for syntax and symbol, after which the Upper Paleolithic creative explosion, the successful dispersal from Africa, and the displacement of the other hominins follow \cite{klein1995,baryosef2002,mellars2006}. Its linguistic variant is the single mutation that yields Merge \cite{berwick2016}. Its evidence is a contrast: anatomically modern humans in the Levant around a hundred thousand years ago who did not persist, and the lineage that left around seventy thousand years ago and did \cite{hershkovitz2018}.

Fisher gives two objections, and both are objections this book has already made in its own terms.

The first is metabolic. Aiello and Wheeler's expensive-tissue hypothesis puts the running cost of the human brain at a fifth to a quarter of resting metabolism, bought by shrinking the gut \cite{aiello1995}. A lone individual born with such an organ, in a foraging economy at the margin, is not advantaged. In the Turn that individual is Lemma~\ref{sci:lem:exposure}: a mortal computation with a strictly positive burn rate and no free names performs no harvest, its stack is monotone decreasing, and it dies in a bounded number of steps. The lemma is about closure, not about brains, but the brain is a burn rate, and a brain whose use requires a counterparty that does not yet exist is a closed term with a large burn rate. The expensive-tissue objection is the exposure lemma with a number attached.

The second is structural. Fisher asks who the first talking ape talks to, and the answer in the calculus is nobody, and the reason is the only rule the calculus has. A send with no complementary receive in its namespace performs no step; a receive with no send performs no step; the \textsc{comm} rule fires on a pair or not at all. A capacity that pays only when a counterparty has it cannot spread by appearing in one individual, because in one individual it is a half of a redex. This is not an analogy. It is the shape of the primitive.

Now the part that is the Prestige's business. Section~\ref{sec:the-move} located the one place in the Turn where something was assumed to exist at no cost: the encoding $\lceil X \rceil$ in Chapter~\ref{sec:isolation}, written with a pair of angle brackets and never constructed, through which the entire outside world enters the theory.

\begin{observation}[The monolith is the angle brackets]
\label{st:obs:monolith}
The monolith story and Proposition~\ref{prop:isolation-closure} make the same move. Each posits a correspondence between a mind and a world that appears in a single agent, at a single moment, with no counterparty and no price: in the one case a mutation that lets an ape think about things, in the other an encoding that lets a term interact with an entity outside the theory. Each is correct about everything downstream of the posit and silent about the posit. The objections to the monolith --- no individual can pay for it, no individual can use it alone --- are the objections Observation~\ref{obs:social} makes to the encoding: reference is a property of a composite of at least two agents, and no single agent's symbols fix it.
\end{observation}

\noindent
i record this not to score a point against a popular book but because it tells me something about the error. The monolith is the first origin story one learns, and the angle brackets are the first thing one writes when one wants to get on with the mathematics. They are both the move of a mind in a hurry, and the three stories that follow are each a way of slowing down and asking who paid.

What the monolith's evidence actually shows is worth keeping. The failed and successful dispersals are real, and the gap between them is real \cite{hershkovitz2018,harvati2019}. The fourth story will read that gap as a density threshold rather than a mutation, and the reading is better, but the datum belongs to whoever explains it.

\section{We: the join is the trigger}
\label{st:sec:we}

The second story is Tomasello's \cite{tomasello2005,tomasello2008,tomasello2019,tomasello2025}. Its pressure is the drying of the East African Rift, after which the lone foraging that had worked in forest stopped working on open ground. Its trigger is a specific cognitive schema that human infants display at around nine months and no great ape displays at any age: a joint goal held with another agent while playing a distinct role in it, joint attention to a thing while seeing it from a distinct perspective \cite{carpenter1998}. The comparative evidence is the Primate Cognition Test Battery, in which two-and-a-half-year-old children match chimpanzees and orangutans on physical cognition and exceed them only on social cognition \cite{herrmann2007}. Tomasello's evolutionary scenario runs in two steps: obligate collaborative foraging produces joint intentionality, and then life in larger groups with competing out-groups produces collective intentionality, with its norms, its institutions, and its ``we'' \cite{tomasello2012}. What intelligence is, on this story: shared agency, a new kind of agent made of agents, in the sense in which an organism is a new kind of agent made of cells \cite{tomasello2022}.

Here is the same story in the Turn.

\begin{observation}[Shared intentionality is a join]
\label{st:obs:join}
A joint goal with distinct roles, and joint attention with distinct perspectives, is the shape of a multi-party rendezvous: one cut, two sides, a continuation that neither side can fire alone. In the calculus it is a join pattern, $\mathtt{for}(x \leftarrow a;\; y \leftarrow b)\,P$, whose body runs only when both channels have been supplied and whose two binders are read from two different places. Chimpanzee cooperation, on Tomasello's description, is two independent redexes that happen to reduce near each other; human cooperation is one redex with two parties. The difference Herrmann and colleagues measure --- no gap on physical cognition, a gap on social cognition --- is the difference between the term language, which the two species share, and the interaction rules, which they do not.
\end{observation}

\noindent
The pressure is the funding gate. When a harvest can no longer be completed by one party, the gate $\chi$ of Chapter~\ref{ch:grading} closes on every single-party rewrite that used to feed, and the only rewrites that remain funded are the joins. That is what ``could not eat alone'' means in a theory where eating is a receipt. Selection for the best coordinators is then not a hypothesis added to the model; it is what the gate does to a population in which some members can complete joins and others cannot.

Two further features of the story are already in Chapter~\ref{ch:com}, and were put there before i had read Tomasello.

The first is Observation~\ref{com:obs:amort}, cooperation is amortized perception: there are two routes to any predicate about an environment, reading it or being told it by a learner who has already paid to read it, and the second is cheaper whenever the reading cost exceeds the rendezvous cost. That is the economics of a shared goal stated as a cost relocation. The second is Tomasello's two steps, joint and then collective, which are the two levels of Chapter~\ref{ch:com}: the dyad composed by a shared name, and the population in which the learner is a distribution and the correspondence is carried by no member alone (Section~\ref{sec:social}).

Where the story and the calculus disagree is instructive. Proposition~\ref{com:prop:nomet} says there is no metabolic mutualism: every transfer of tokens is zero-sum, so cooperation cannot be a token exchange that leaves both parties richer. Tomasello's foragers are richer for cooperating. The reconciliation is Observation~\ref{com:obs:two}, tokens are conserved and formulae are not: the surplus of a shared hunt is not in the division of the carcass, which is zero-sum, but in the fact that a predicate about where the antelope is, once paid for by one, is held by both at no further cost. The ``we'' is positive-sum exactly where it is epistemic, and the calculus says this more sharply than the anthropology does.

A caution about dates. Fisher pins the start of the cooperation story to eight hundred thousand years ago, on the first controlled fire at Gesher Benot Ya\textquoteleft aqov and the first big-game hunting \cite{goreninbar2004,thieme1997}; Tomasello's own scenario puts the interdependence step at around four hundred thousand years and the collective step at around a hundred and fifty thousand \cite{tomasello2012}. The calculus has no opinion about dates. It has an opinion about order: the join must be in use before anything in the third or fourth story can begin, because both of those are things a population does \emph{to its own members}, and that requires the population to be able to act as one. This is the sense in which Fisher is right that story two is prerequisite to the others, and it is a structural rather than a historical claim.

\section{Schemers: the coalition can starve you}
\label{st:sec:schemers}

The third story is really two literatures that Fisher, reasonably, fuses. i will keep them apart because the calculus treats them differently.

\subsection{Self-domestication}

The first literature is self-domestication \cite{wrangham2019,wrangham2018,hare2017,hare2012}. Its evidence begins with a morphological signal: from around three hundred thousand years ago sapiens skulls gracilize, with shorter faces, smaller teeth, reduced brow ridges, and later a modest reduction in brain volume \cite{cieri2014,desilva2021}. That package is the domestication syndrome, reproduced experimentally in Trut's Siberian farm-fox lineage by selection on tameness alone \cite{trut1999,dugatkin2017}, and given a developmental mechanism by Wilkins, Wrangham and Fitch as mild hypofunction of the neural crest, which supplies the adrenal medulla and also cartilage, pigment cells and parts of the face \cite{wilkins2014}. The genomic comparison of sapiens with Neanderthals finds overlap between the genes under selection in us and those under selection in domesticates relative to their wild relatives \cite{theofanopoulou2017}. The inference is Wrangham's: somebody domesticated us, and it was us. Once coalitions of subordinate males could conspire, they could kill a tyrannical alpha, and the universal sanction in forager societies against the would-be despot is execution \cite{boehm1999,boehm2012}. That inverted the primate selection gradient. Reactive aggression became lethal to its bearer.

The calculus has a place for each piece.

\begin{observation}[Execution is forcing by withdrawal]
\label{st:obs:coalition}
A coalition that can kill a member is, in the vocabulary of Chapter~\ref{ch:chs}, a coalition that can force an outcome: a set of agents whose joint strategy decides a proposition whatever the rest do. In a conserved economy the only outcome a coalition needs to force is a debit the target cannot cover. Death in the Turn is a stack reaching zero (Lemma~\ref{sci:lem:exposure}), and a coalition that can refuse every receipt on a member's metabolic names has forced it. Execution requires no new primitive. It requires that the join of Section~\ref{st:sec:we} be available to a subset against a member, which is the same construction aimed inward.
\end{observation}

\noindent
This is also where the Pledge's warning about ambient authority (Section~\ref{sec:pledge-ii}) becomes an evolutionary pressure. The question of who may spend is the question of who may cut whom off, and a population in which that question has an answer is one in which being on the wrong side of the answer is fatal.

The second piece is where self-domestication lives in the formalism, and i think the placement is the most useful thing in this section.

\begin{observation}[Self-domestication changes the resolver, not the terms]
\label{st:obs:resolver}
Selection against reactive aggression does not change what the members of a population can do; it changes which of their possible interactions the population permits to proceed. In Chapter~\ref{ch:grading} that distinction is exact: the term language is fixed, the race among candidate redexes is decided by a resolver, and Principle~\ref{gw:prin:aperture} says the resolver is a parameter that no context reads. A self-domesticating population is one whose resolution algebra (Definition~\ref{gw:def:algebra}) has been revised so that candidates in which one member seizes another's resources are graded toward zero. The animal is the same term. The ecology has changed what the clause returns.
\end{observation}

\noindent
Two things follow, one reassuring and one not. The reassuring one is support adequacy (Corollary~\ref{gw:cor:adequacy}): whatever a revised resolver returns is supported on behaviors the unweighted calculus already had, so a domesticated population is not a different species of computation, which is what the morphological continuity says. The unreassuring one is that the foxes' brains shrank, and under this reading one expects it: if selection acts on the resolver, the pressure on the term language is relaxed, and whatever in the term language was expensive and no longer load-bearing is free to go. The reading predicts gracilization as a relaxation, not as an adaptation, which is what the byproduct premise would want.

The honest objections. The brain-shrinkage datum is disputed in its timing; DeSilva and colleagues place a measurable decrease only around three thousand years ago, and Villmoare and Grabowski contest even that \cite{desilva2021,villmoare2022}. Lord and colleagues argue that the farm foxes were already farm-bred when Trut's experiment began and that the domestication syndrome is not the tidy package the story needs \cite{lord2020}. Neither objection touches Observation~\ref{st:obs:resolver}, which is about where the construction sits, but both touch how much weight the human case can bear.

\subsection{Machiavellian intelligence}

The second literature is older and separate: Humphrey's social function of intellect, the Machiavellian intelligence program of Byrne and Whiten, and Dunbar's social brain hypothesis correlating neocortex ratio with group size \cite{humphrey1976,byrne1988,dunbar1998}. Its claim is that the driver of primate brain expansion was the complexity of conspecifics, not of ecology. Its best current challenge is DeCasien and colleagues' finding that across primates brain size is predicted by diet and not by sociality \cite{decasien2017}; Wrangham's tyranny of cousins, which he takes from Gellner \cite{gellner1988}, is the form in which the social pressure becomes lethal rather than merely competitive, and so survives that challenge in the human case even if it fails in the primate one.

The Turn already prices this.

\begin{remark}[Modeling a modeler is quoting a quoter]
\label{st:rem:machiavelli}
To navigate a coalition one must model what the other members will do, which requires modeling what they believe one will do. In the calculus a model of another learner is a quotation, $@P$, held and inspected; a model of another learner's model of oneself is a quotation of a quotation; and Section~\ref{sec:pledge-iii} already described the exponential explosion this produces and the aggressive compression that keeps it affordable. Reflection is what makes the regress expressible in one move. Cost accounting is what bounds it. The Machiavellian brain is a reflective learner whose budget for nested quotation has been raised by selection, and the limit on political depth in any real coalition is the limit of Section~\ref{sec:kant-affordable}: modal depth is a price.
\end{remark}

\noindent
And this is where the third story meets the chapter before this one. Chapter~\ref{ch:lie} argued that meaning is the receiver's forced expense, and that the first symbol was a lie because a lie is the first signal whose reception costs something to check (Observation~\ref{lie:obs:thesis}). A coalition is an audience of receivers who can impose that expense on the sender. The tyranny of cousins is the ecological condition under which Benedikt's mechanism has teeth: a population that can execute is a population before which a false signal is expensive, and the Machiavellian brain is the organ that computes, under that threat, what a signal will cost.

What intelligence is, on the third story, is social cognition. In the ledger it is the depth of reflection a learner can fund under a resolver that its own population controls. Fisher's observation that this explains tribalism, conformism and conspiracy belief has a formal shadow: those are behaviors that lower the price of staying inside a coalition, and a learner that optimizes the ledger rather than the truth will pay for them.

\section{Hivemind: culture is quoted code, and density is compactness}
\label{st:sec:hivemind}

The fourth story is Henrich's, with Boyd and Richerson behind him \cite{henrich2016,boyd1985,richerson2005,boyd2011}. It begins with Jebel Irhoud, which puts anatomically modern sapiens at around three hundred thousand years \cite{hublin2017}, and with the consequence that for a quarter of a million years the archaeological record of modern behavior flickers: present here, absent for tens of millennia, present again elsewhere \cite{mcbrearty2000}. Henrich's account is that what evolved was high-fidelity social learning --- imitate the successful, defer to prestige, over-copy opaque procedures --- and that its payoff depends on two things outside the learner. The first is the rate of environmental change, which must be fast enough that genes cannot track it and slow enough that a parent's knowledge is still useful: Potts's variability selection, with the Olorgesailie record of Middle Stone Age technology and long-distance obsidian exchange around three hundred and twenty thousand years ago as the type case \cite{potts1998,potts2018}. The second is population size and connectedness. Tasmania's toolkit lost complexity after isolation; Oceanic island toolkits scale with population; density thresholds predict when modern behavior appears in the Late Pleistocene record \cite{henrich2004,kline2010,powell2009}; transmission-chain experiments reproduce the effect in the laboratory \cite{derex2013}. The result Henrich calls the collective brain \cite{muthukrishna2016}, with the cultural brain hypothesis as the formal model of why brain size, sociality and social learning co-vary \cite{muthukrishna2018}. What intelligence is, on this story: a property of a connected population, with individual cleverness a weak proxy for it.

Of the three surviving stories this is the one the Turn already contains most completely, and i was not expecting that.

\begin{observation}[Culture is quotation, and quotation is non-rival]
\label{st:obs:nonrival}
Henrich's ``second inheritance system'' is Observation~\ref{com:obs:two}: tokens are conserved and formulae are not. A practice transmitted by imitation is a quotation $@P$ held by the imitator; holding it costs nothing while held, acquiring it is grade-three access, the cheapest on the ladder, and one learner's acquiring it deprives no other. The cassava procedure is a formula nobody in the chain can read and everybody can hold. That it is \emph{code} and not merely a disposition is what reflection adds: an inherited quotation can be dropped, run, edited and re-quoted, which is what the ratchet requires and what a fixed rewrite rule cannot do. The reproduction note's crossover on quoted code is the mechanism of cumulative culture stated at the level of the primitive.
\end{observation}

\begin{observation}[The ratchet is assembly index at copy number]
\label{st:obs:ratchet}
Cumulative culture is Chapter~\ref{sec:assembly-depth}'s two axes read as one process. The assembly index of a practice is its depth, the number of joining steps that had to be paid to build it; copy number is how many learners hold it. Theorem~\ref{thm:ai-bound} says depth bounds what can be told apart, and Remark~\ref{rmk:depth-width} says depth is exponentially fragile and width linearly cheap, so a lineage under scarcity loses levels before it loses population. Henrich's Tasmania is that remark run as history: a population below the width needed to sustain its depth loses the deep practices first. Proposition~\ref{scm:prop:frontier} makes the pair a frontier in the $(\AI,\CN)$ plane whose slope is set by medium reliability, and the frontier is the collective brain's budget line.
\end{observation}

\begin{observation}[Density is the compactness condition]
\label{st:obs:density}
The bottleneck Henrich calls density is the condition the book has approached twice from the other side. Section~\ref{sec:pledge-vii} argued that a population can agree only up to the scale at which messages still arrive intact, and Conjecture~\ref{conj:degradation} states it: beyond a scale $N_\varepsilon$ no agent communicates coherently with the whole population. Chapter~\ref{sec:topology} then named the condition under which a population's hypotheses converge --- compactness of the logical topology --- and Remark~\ref{rmk:consensus-missing} located what compactness alone does not supply, which is termination, and found it in the budget. Henrich's claim is the same condition from below: a practice accumulates only within a population connected enough that a revision made anywhere reaches everywhere before the budget that made it runs out. Powell, Shennan and Thomas's thresholds and Section~\ref{sec:pledge-vii}'s clusters are the two ends of one rope, and the flickering record of the Middle Stone Age is what one expects of populations that cross and re-cross the threshold.
\end{observation}

\noindent
The variability-selection pressure also lands on a specific object, and it is the one the whole of this iteration of the book has been built around.

\begin{observation}[Variability selection is the plasticity hypothesis]
\label{st:obs:variability}
Potts's window --- environmental change too fast for genes, slow enough for parents --- is the hypothesis of Proposition~\ref{gw:prop:plastic} stated ecologically. A gene is a rewrite rule: adaptation by genes is revision of the theory, which is slow because it is revision of the rule set. A learned practice is a datum in the term that a clause reads: adaptation by culture is revision of what the clause is evaluated against, which is fast because the clause is re-evaluated at every candidate and no rule changes (Remark~\ref{gw:rmk:state-is-term}). The proposition's hypothesis is that the trace dependence be expressible in the current term, and Henrich's two-sided window is the condition under which it is worth paying to make it so: fast enough that rule revision loses, slow enough that the term still carries usable information from the previous generation. When the environment is slower than that, encode it in the rules; when it is faster, nothing helps.
\end{observation}

\noindent
The two caveats that belong here. Vaesen and colleagues argue that population size does not explain past changes in cultural complexity in the archaeological record as cleanly as the model wants \cite{vaesen2016}; and the Eurasian east-west axis that Henrich takes from Diamond \cite{diamond1997} is a geographical claim about connectivity on which the calculus is silent except to say that connectivity is the variable and not population count as such, which is also what Observation~\ref{st:obs:density} says.

\section{Three pressures, one ledger}
\label{st:sec:ledger}

Fisher's closing judgment is that the three surviving stories crisscross and overlap and are probably all partially true, and that story two is needed to get to either of the others. The ledger lets me say that more precisely, and say where it is wrong to leave it there.

The three stories are three constraints on one account, and they act on different parts of it.

\begin{itemize}[leftmargin=2em]
\item \textbf{We} is a constraint on the \emph{funding gate}: no single-party rewrite feeds, so only joins are funded. It is a precondition because the other two constraints are things a population does, and a population that cannot join cannot do anything.
\item \textbf{Schemers} is a constraint on the \emph{resolver}: the population grades candidate interactions, and the grading can starve a member. It acts inward, on who may spend.
\item \textbf{Hivemind} is a constraint on \emph{transmission}: the population's depth is bounded by its width and its connectivity, and its rate of adaptation by whether the environment lets a clause do the work of a rule. It acts outward, on who may copy.
\end{itemize}

\noindent
Chapter~\ref{ch:com} already says that the typed overlap between learners is a vector and not a scalar (Section~\ref{com:sec:typed}), and the three stories are three coordinates of that vector under pressure. That is why ``which story is true'' is the wrong question. A learner under all three constraints is a different object from a learner under any one, and the human record is the trace of the composite.

The place where the stories genuinely conflict is also visible in the vector. Self-domestication lowers behavioral variance: it grades the hothead toward zero, and in the foxes it narrowed the distribution of temperament. The cultural ratchet needs variance: innovation is a rare candidate that the population must sometimes let fire. A resolver tuned hard enough to suppress the aggressive outlier will also suppress the inventive one, and a population whose resolver is loose enough to keep its inventors will keep some of its despots. i do not think this tension is resolved in the literature, and the ledger does not resolve it; it locates it, in the grading of the overlap vector, and says that it is a trade-off and not a contradiction.

\begin{conjecture}[Flicker, then ignition]
\label{st:conj:flicker}
Consider a population of mortal scientists of Chapter~\ref{ch:sci}, composed as in Chapter~\ref{ch:com}, under a funding gate that admits only joins on the metabolic names, a resolver under the population's own control that can grade any member's receipts to zero, and a copy channel on quoted practices subject to Conjecture~\ref{conj:degradation}. Vary the population's size and connectivity across the compactness threshold of Chapter~\ref{sec:topology} and the environment's rate of change across the window of Observation~\ref{st:obs:variability}. Then the depth of the deepest practice held by any member exhibits the profile the Middle Stone Age record shows: long intervals of loss and recovery below the threshold, and monotone accumulation above it, with the transition controlled by connectivity rather than by any change in the members.
\end{conjecture}

\noindent
This is a statement about a simulation nobody has run, and i state it as a target because every piece of it is already an object in the book and the composite is the first thing i would build. What would falsify it is a population above the threshold that still flickers, or one below it that accumulates; what would falsify the whole reading is a lineage that accumulated depth while losing width, which Remark~\ref{rmk:depth-width} forbids.

\section{What this changes about the thing we are building}
\label{st:sec:ai}

Fisher ends by asking, for each story, whether it makes one more or less confident that an artificial system could follow the same route. His answers are that the schemers' story is pessimistic for machines, which do not live under coalitions that can kill them; that the ``we'' story is cautiously optimistic, since agents increasingly coordinate; and that the hivemind story is the one that should worry or encourage us most, since training on an ever-larger corpus looks a great deal like a collective brain.

The ledger sharpens all three, and the sharpening is not comfortable.

On the first: a system that cannot be cut off by the population it serves has no resolver under that population's control, and Observation~\ref{st:obs:resolver} says that the social intelligence of the third story is exactly what such a resolver selects for. The absence is not a gap in the system's training. It is the absence of the construction. Section~\ref{sec:blind-spot-returned} already said that such a system will not refer on its own; this adds that it will not scheme on its own either, and that a system which appears to is reproducing the surface of a behavior whose selecting pressure it has never been under. Whether that is a relief depends on what one wanted.

On the second: the join is cheap to build and its selecting pressure is cheap to impose --- a funding gate that admits only multi-party rewrites is a line in a resolver --- and nothing in the Turn says that a population of artificial learners under that gate would fail to produce the first story's trigger. What it says is that the trigger is the join and not the brain, which is a different thing to build from what is currently being built.

On the third: the hivemind reading is the one the ledger most strongly endorses and most strongly qualifies. A corpus is a store of quotations and training on it is grade-three access at scale; that is a collective brain in the sense of Observation~\ref{st:obs:nonrival}. But Observation~\ref{st:obs:ratchet} says the collective brain's depth is bounded by its width \emph{at the time the depth was built}, and a system that reads a corpus inherits the depth of the population that wrote it without inheriting the width that sustains it. Henrich's Tasmania is the warning: depth without the population to maintain it is depth about to be lost. A system that is a collective brain by inheritance and not by membership is a Tasmania, and Chapter~\ref{ch:psy} already named what that loss looks like from the inside (Principle~\ref{psy:prin:lapse}).

\section{What is claimed, and what is not}
\label{st:sec:claims}

Claimed: that Fisher's pressure--trigger--class schema is the schema of the Turn (Principle~\ref{st:prin:schema}); that the byproduct premise the three surviving stories share is the ledger premise of this book (Observation~\ref{st:obs:byproduct}); that the monolith and the encoding of Chapter~\ref{sec:isolation} are one move (Observation~\ref{st:obs:monolith}); and that each of the three surviving stories names a constraint already present in the Turn as a distinct construction --- the join under the funding gate, the resolver under the population's control, and the copy channel under the compactness and plasticity conditions (Observations~\ref{st:obs:join}, \ref{st:obs:resolver}, \ref{st:obs:nonrival}--\ref{st:obs:variability}).

Not claimed: that the anthropology confirms the calculus or the calculus the anthropology. Convergence is not confirmation, and the same discipline Remark~\ref{suf:rem:solms} applied to Solms applies here. Not claimed: Conjecture~\ref{st:conj:flicker}, which has not been run. Not claimed: that the dates are right, which the calculus cannot assess, or that the empirical disputes named in each section are settled, which they are not.

Also not claimed, and worth saying because the author of the Pledge chapter is the same person: that any of this tells one which story to prefer. i prefer the schemers, as Fisher does, and i have given the reason in the Pledge and said why it is a bad one. The ledger's answer is that preference is misplaced, because the stories are coordinates and not candidates. The human record is what one composite account looked like under three pressures at once. The interesting question is not which pressure but what the books looked like, and the Turn is an attempt to write books of that kind.
